\section{Generated Art Anova - For all components touched for all commits} 
An \textbf{Aligned Rank Transform (ART)} Anova table for all components touched on average. 
The statistical significance of architectural exploration is evaluated using a dual-framework modeling methodology. A standard parametric ordinary least squares (\textbf{OLS}) ANOVA is presented alongside a robust non-parametric ANOVA. This side-by-side presentation explicitly demonstrates where extreme right-skew and outlier-induced variance in repository mining data mask true behavioral effects in traditional linear models, which the rank-transformed framework uncovers.

\textbf{Explanation of Combined Table Columns:}
\begin{itemize}
    \item \textbf{Component / Commits ($n$) / Devs ($N$):} Define the experimental scope, tracking the architectural tier (packages, classes, or methods), the longitudinal lifecycle snapshot threshold, and the active developer population size.
    \item \textbf{Factor:} The independent categorical variables or interactions being tested for statistical significance.
    
    \item \textbf{OLS df:} Parametric numerator degrees of freedom based on factor grouping configurations.
    \item \textbf{OLS F-Stat:} The traditional variance ratio metric calculated from raw values.
    \item \textbf{OLS $P$-Value:} Parametric significance probability. Note how heteroscedastic noise frequently inflates these values toward non-significance ($p > 0.05$).
    \item \textbf{OLS $\eta^2$ (Eta-Squared):} The total proportion of raw variance explained exclusively by this standalone factor.
    
    \item \textbf{ART df (Resid):} Formatted as $df_{\mathrm{effect}}~(df_{\mathrm{residual}})$. Because the ART methodology conducts independent alignment procedures for each factor before ranking, it assigns the exact denominator error degrees of freedom dynamically alongside its numerator.
    \item \textbf{ART F-Stat:} The robust variance ratio calculated over the aligned, ranked data layers.
    \item \textbf{ART $P$-Value:} Non-parametric significance probability, providing the authoritative test for heavy-tailed software metrics.
    \item \textbf{ART Part. $\eta^2$ (Partial Eta-Squared):} The effect size over ranked data layers, indicating the proportion of variance explained by this specific factor after controlling for alternate main effects.
\end{itemize}

\textbf{Explanation of Fixed Table Rows:}
\begin{itemize}
    \item \textbf{Role (F/LJ):} Isolates the \textit{Main Role Effect} to evaluate whether the timing of project entry (Founder vs. Later Joiner) significantly alters exploration breadth on its own.
    \item \textbf{Category (T/M/S or M/S):} Isolates the \textit{Main Category Effect} to determine whether a contributor's eventual longevity footprint (Transient, Moderate, or Sustained) dictates early behavior.
    \item \textbf{Interaction:} Evaluates the \textit{Interaction Effect} ($\mathrm{Role} \times \mathrm{Category}$). A significant interaction indicates that the impact of entry role is not uniform, but behaves completely differently depending on the eventual tenure category of the contributor.
\end{itemize}
\begin{table}[H]
\centering
\setlength{\tabcolsep}{2pt}
\caption{ART Anova results for all components, number of commits and categories of developers}
\resizebox{\textwidth}{!}{% 
\begin{tabular}{l r r l ccc r ccc r}
\toprule
\textbf{Component} & \makecell[t]{\textbf{First $n$}\\\textbf{commits}} & \makecell[t]{\textbf{Number of}\\\textbf{Developers}} & \textbf{Factor} & \multicolumn{4}{c}{\textbf{OLS}} & \multicolumn{4}{c}{\textbf{ART}} \\ 
\cmidrule(lr){5-8} \cmidrule(lr){9-12}
 &  &  &  & \textbf{df} & \textbf{F-Stat} & \textbf{$P$ Value} & \textbf{$\eta^2$} & \textbf{df (Resid)} & \textbf{F-Stat} & \textbf{$P$ Value} & \makecell[t]{\textbf{Part.}\\\textbf{$\eta^2$}} \\ 
\midrule
packages & 1 & 16434 & Role (F/LJ) & 1 & 0.00 & 0.9606 & 0.0000 & 1 (16428) & 265.55 & 0.0000 & 0.0159 \\ 
packages & 1 & 16434 & Category (T/M/S) & 2 & 32.56 & 0.0000 & 0.0039 & 2 (16428) & 801.56 & 0.0000 & 0.0889 \\ 
packages & 1 & 16434 & Interaction & 2 & 0.60 & 0.5510 & 0.0001 & 2 (16428) & 132.65 & 0.0000 & 0.0159 \\ 
packages & 2 & 3085 & Role (F/LJ) & 1 & 0.00 & 0.9820 & 0.0000 & 1 (3079) & 6.11 & 0.0135 & 0.0020 \\ 
packages & 2 & 3085 & Category (M/S) & 2 & 0.32 & 0.7261 & 0.0002 & 2 (3079) & 8.25 & 0.0003 & 0.0053 \\ 
packages & 2 & 3085 & Interaction & 2 & 0.12 & 0.8843 & 0.0001 & 2 (3079) & 0.43 & 0.6534 & 0.0003 \\ 
packages & 5 & 1362 & Role (F/LJ) & 1 & 0.00 & 0.9828 & 0.0000 & 1 (1356) & 0.18 & 0.6679 & 0.0001 \\ 
packages & 5 & 1362 & Category (M/S) & 2 & 0.20 & 0.8187 & 0.0003 & 2 (1356) & 4.70 & 0.0093 & 0.0069 \\ 
packages & 5 & 1362 & Interaction & 2 & 0.12 & 0.8879 & 0.0002 & 2 (1356) & 0.02 & 0.9789 & 0.0000 \\ 
packages & 10 & 857 & Role (F/LJ) & 1 & 0.03 & 0.8676 & 0.0000 & 1 (853) & 0.25 & 0.6192 & 0.0003 \\ 
packages & 10 & 857 & Category (M/S) & 1 & 0.41 & 0.5234 & 0.0005 & 1 (853) & 0.42 & 0.5167 & 0.0005 \\ 
packages & 10 & 857 & Interaction & 1 & 0.16 & 0.6929 & 0.0002 & 1 (853) & 1.32 & 0.2507 & 0.0015 \\ 
classes & 1 & 16434 & Role (F/LJ) & 1 & 47.88 & 0.0000 & 0.0029 & 1 (16428) & 463.03 & 0.0000 & 0.0274 \\ 
classes & 1 & 16434 & Category (T/M/S) & 2 & 34.71 & 0.0000 & 0.0042 & 2 (16428) & 912.14 & 0.0000 & 0.0999 \\ 
classes & 1 & 16434 & Interaction & 2 & 23.37 & 0.0000 & 0.0028 & 2 (16428) & 305.46 & 0.0000 & 0.0359 \\ 
classes & 2 & 3085 & Role (F/LJ) & 1 & 9.04 & 0.0027 & 0.0029 & 1 (3079) & 216.75 & 0.0000 & 0.0658 \\ 
classes & 2 & 3085 & Category (M/S) & 2 & 1.91 & 0.1484 & 0.0012 & 2 (3079) & 24.78 & 0.0000 & 0.0158 \\ 
classes & 2 & 3085 & Interaction & 2 & 3.96 & 0.0192 & 0.0025 & 2 (3079) & 9.47 & 0.0001 & 0.0061 \\ 
classes & 5 & 1362 & Role (F/LJ) & 1 & 4.76 & 0.0293 & 0.0034 & 1 (1356) & 13.56 & 0.0002 & 0.0099 \\ 
classes & 5 & 1362 & Category (M/S) & 2 & 0.94 & 0.3914 & 0.0014 & 2 (1356) & 17.91 & 0.0000 & 0.0257 \\ 
classes & 5 & 1362 & Interaction & 2 & 2.19 & 0.1119 & 0.0032 & 2 (1356) & 3.97 & 0.0191 & 0.0058 \\ 
classes & 10 & 857 & Role (F/LJ) & 1 & 3.15 & 0.0765 & 0.0036 & 1 (853) & 24.98 & 0.0000 & 0.0285 \\ 
classes & 10 & 857 & Category (M/S) & 1 & 1.27 & 0.2592 & 0.0015 & 1 (853) & 54.96 & 0.0000 & 0.0605 \\ 
classes & 10 & 857 & Interaction & 1 & 2.71 & 0.1000 & 0.0031 & 1 (853) & 11.56 & 0.0007 & 0.0134 \\ 
methods & 1 & 16434 & Role (F/LJ) & 1 & 12.69 & 0.0004 & 0.0008 & 1 (16428) & 382.32 & 0.0000 & 0.0227 \\ 
methods & 1 & 16434 & Category (T/M/S) & 2 & 10.01 & 0.0000 & 0.0012 & 2 (16428) & 801.24 & 0.0000 & 0.0889 \\ 
methods & 1 & 16434 & Interaction & 2 & 9.24 & 0.0001 & 0.0011 & 2 (16428) & 374.01 & 0.0000 & 0.0436 \\ 
methods & 2 & 3085 & Role (F/LJ) & 1 & 2.38 & 0.1227 & 0.0008 & 1 (3079) & 255.19 & 0.0000 & 0.0765 \\ 
methods & 2 & 3085 & Category (M/S) & 2 & 0.50 & 0.6076 & 0.0003 & 2 (3079) & 38.72 & 0.0000 & 0.0245 \\ 
methods & 2 & 3085 & Interaction & 2 & 1.76 & 0.1727 & 0.0011 & 2 (3079) & 24.86 & 0.0000 & 0.0159 \\ 
methods & 5 & 1362 & Role (F/LJ) & 1 & 1.16 & 0.2815 & 0.0009 & 1 (1356) & 43.86 & 0.0000 & 0.0313 \\ 
methods & 5 & 1362 & Category (M/S) & 2 & 0.20 & 0.8199 & 0.0003 & 2 (1356) & 26.99 & 0.0000 & 0.0383 \\ 
methods & 5 & 1362 & Interaction & 2 & 0.87 & 0.4212 & 0.0013 & 2 (1356) & 6.77 & 0.0012 & 0.0099 \\ 
methods & 10 & 857 & Role (F/LJ) & 1 & 0.78 & 0.3789 & 0.0009 & 1 (853) & 78.52 & 0.0000 & 0.0843 \\ 
methods & 10 & 857 & Category (M/S) & 1 & 0.27 & 0.6054 & 0.0003 & 1 (853) & 64.36 & 0.0000 & 0.0702 \\ 
methods & 10 & 857 & Interaction & 1 & 1.09 & 0.2968 & 0.0013 & 1 (853) & 11.05 & 0.0009 & 0.0128 \\ 
\bottomrule
\end{tabular}%
} 
\label{tab:codebase_graph_latex_repository_graph_component_time_series_generated_ART_anova} 
\end{table}
\newpage 


\subsubsection*{Statistical Model}
The differences between developer categories are assessed using a Two-Way ANOVA 
with interaction. The value for any individual developer metric ($Y_{ijk}$) 
is modeled as:

\begin{equation}
    Y_{ijk} = \mu + \alpha_i + \beta_j + (\alpha\beta)_{ij} + \epsilon_{ijk}
\end{equation}

Where:
\begin{itemize}
    \item $\mu$ is the grand mean of the touched components.
    \item $\alpha_i$ represents the effect of the \textbf{Role} (Founder vs. Later Joiner).
    \item $\beta_j$ represents the effect of the \textbf{Tenure} (Transient, Moderate, Sustained).
    \item $(\alpha\beta)_{ij}$ is the interaction effect between Role and Tenure.
    \item $\epsilon_{ijk}$ is the residual error.
\end{itemize}

The $F$-statistic for each factor is calculated by:
\begin{equation}
    F = \frac{SS_{factor} / df_{factor}}{SS_{residual} / df_{residual}}
\end{equation}

\newpage 
